\chapter*{第 1 章 绪论}
\addcontentsline{toc}{chapter}{第 1 章 绪论}
\markboth{第 1 章 绪论}{}

\begin{tishi}
本章为绪论，818 初试\textbf{不考大题}，只出选择、填空、判断、简答中 1--3 题。
做题要求：概念题\textbf{写出完整要点}（不是"略"），计算题必须写出公式、代入、单位、结果四步。
\end{tishi}

\begin{ti}
\sloppy
\emergencystretch=2em

\item 何谓流体的连续介质模型？为了研究流体机械运动规律，说明引入连续介质模型的必要性。\xstarfull{4}\yiju{E4 2022 年 818 单选第 1 题考"连续介质假设的含义"；E1 slide33"第一章绪论＝重点掌握基础概念"；E5 题库问答题第 2 题}\nandu{易}

\item 试述牛顿内摩擦定律？产生摩擦力的根本原因是什么？\xstarfull{2}\yiju{教材式 1.9、1.10（E6）；E4 2026 年回忆版计算题第 1 题为牛顿流体切应力问题；E5 题库问答题第 1 题"流动性与黏滞性"}\nandu{易}

\item 液体和气体的黏性随温度的升高或降低发生变化，变化趋势是否相同？为什么？\xstarfull{4}\yiju{E4 2026 年回忆版简答考"层流与雷诺数关系"、概念涉及理想流体；E2"第 1--3 章偏概念，掌握基本公式与解释即可"；教材表 1.1 与 1.3.3 节结论（E6）}\nandu{易}

\item 作用在流体上的外力有哪两类？是如何定义的？\xstarfull{3}\yiju{E5 题库选择第 3 题"体积力有势"；E6 教材 1.4 节（第 2 章静压强特性的直接来源）}\nandu{易}

\item 如习题 1.5 图所示，有一 $0.8\ \mathrm{m}\times0.2\ \mathrm{m}$ 的平板在油面上做水平运动，已知运动速度 $U=1\ \mathrm{m/s}$，平板与固定边界的距离 $\delta=1\ \mathrm{mm}$，油的动力黏度 $\mu=1.15\ \mathrm{Pa\cdot s}$，由平板所带动的油的速度成直线分布，试求平板所受的阻力。
\tuyi{习题 1.5 图：下部为\textbf{固定边界}，其上为\textbf{厚度 $\delta=1\ \mathrm{mm}$ 的油层}；平板平放在油面上（平板本身很薄，厚度忽略），平板以 $U=1\ \mathrm{m/s}$ 水平运动，油层内速度由固定边界处的 $0$ 直线增大到平板处的 $U$。平板的湿面积＝$0.8\ \mathrm{m}\times0.2\ \mathrm{m}$。}
\xstarfull{2}\yiju{E4 2026 年回忆版计算题第 1 题（牛顿流体、切应力）；E2"第 1--3 章偏概念，掌握基本公式与解释即可"；教材式 1.9、1.10（E6）；无直接对应题号证据，故压至两星}\nandu{易}

\item 如习题 1.6 图所示，在相距 $\delta=40\ \mathrm{mm}$ 的两平行平板间充满动力黏度 $\mu=0.7\ \mathrm{Pa\cdot s}$ 的液体，液体中有一长为 $a=60\ \mathrm{mm}$ 的薄平板以 $U=15\ \mathrm{m/s}$ 的速度水平向右移动。假定平板运动引起液体流动的速度分布是线性分布。求：（1）当 $h=10\ \mathrm{mm}$ 时，薄板单位宽度受到的阻力 $T$；（2）当 $h$ 为多大时，薄板的阻力最小，并计算最小值。
\tuyi{习题 1.6 图：上、下两平行固定平板相距 $\delta=40\ \mathrm{mm}$；其间正中附近有一长为 $a=60\ \mathrm{mm}$ 的薄平板（厚度忽略）水平运动，速度为 $U=15\ \mathrm{m/s}$ 向右；薄平板到上方固定板的距离为 $h$，到下方固定板的距离为 $\delta-h$；上下两油层的速度均由固定板处的 $0$ 线性变化到薄平板处的 $U$，两油层对薄平板同时产生反向的黏滞阻力。}
\xstarfull{2}\yiju{教材式 1.9、1.10 与"速度线性分布 $T=\mu A U/h$"的例题方法（E6）；E2"第 1--3 章偏概念"；E4 历年真题中未直接出现此双液层极值题，故压至两星}\nandu{难}

\item 动力黏度 $\mu=0.172\ \mathrm{Pa\cdot s}$ 的润滑油充满在两个同轴圆柱体的间隙中，外筒固定，内径 $D=12\ \mathrm{cm}$，间隙 $h=0.02\ \mathrm{cm}$，试求：（1）当内筒以速度 $U=1\ \mathrm{m/s}$ 沿轴线方向运动时，内筒表面的切应力 $\tau_1$，如习题 1.7 图（a）所示；（2）当内筒以转速 $n=180\ \mathrm{r/min}$ 旋转时，内筒表面的切应力 $\tau_2$，如习题 1.7 图（b）所示。
\tuyi{习题 1.7 图：两个\textbf{同轴圆柱体}，\textbf{外筒固定}（图中打阴影线），外筒内径 $D=12\ \mathrm{cm}$；内筒与外筒之间的\textbf{环状径向间隙} $h=0.02\ \mathrm{cm}$，间隙内充满动力黏度 $\mu=0.172\ \mathrm{Pa\cdot s}$ 的润滑油。（a）内筒沿\textbf{轴线方向}以速度 $U=1\ \mathrm{m/s}$ 平移，油层速度由内筒表面的 $U$ 沿径向线性降到外筒内壁的 $0$，速度梯度为 $U/h$；（b）内筒绕轴线以转速 $n=180\ \mathrm{r/min}$ \textbf{旋转}，内筒表面的周向线速度 $u=\pi Dn/60$，油层周向速度沿径向线性降到外筒内壁的 $0$，速度梯度为 $u/h$。}
\xstarfull{4}\yiju{E5 题库填空"动力黏度为 $\mu$ 的润滑油充满在两个同轴圆柱体的间隙中，外筒固定，间隙为 $h$，内筒直径为 $d$，当内筒以转速 $n$ 旋转时，内筒表面的切应力等于（\hspace{2cm}）"；E4 2022 年 818 判断题第 2 题（同轴圆柱间隙内筒表面切应力 $\tau=\mu U/h$）；E1 slide33"第一章绪论＝重点掌握基础概念"}\nandu{易}

\item 某油的密度为 $851\ \mathrm{kg/m^3}$，运动黏度为 $3.39\times10^{-6}\ \mathrm{m^2/s}$，求此油的重度 $\gamma$、比容 $u$ 和动力黏度 $\mu$。
\xstarfull{4}\yiju{E4 2015 年试题填空第 1 题（原文即"某油的密度为 851 kg/m$^3$，运动粘度为 $3.39\times10^{-6}$ m$^2$/s，此油的重度 $\gamma=$"）；E2"第 1--3 章偏概念，掌握基本公式与解释即可"；教材式 1.3、1.4、1.11（E6）}\nandu{易}

\item 存放 $4\ \mathrm{m^3}$ 液体的储液罐，当压强增加 $0.5\ \mathrm{MPa}$ 时液体体积减小 $1\ \mathrm{L}$，求该液体的体积模量。
\xstarfull{2}\yiju{教材式 1.5--1.7（体积压缩系数与体积弹性模量的定义，E6）；E2"第 1--3 章偏概念、掌握基本公式与解释即可"；E4 历年真题中未见独立成题，故压至两星}\nandu{易}

\item 为防止水温升高时体积膨胀将水管胀裂，通常在水暖系统顶部设有膨胀水箱，如习题 1.10 图所示，若系统内水的总体积为 $10\ \mathrm{m^3}$，加温前后温差为 $50\ ^{\circ}\mathrm{C}$，在其温度范围内水的体积膨胀系数 $\alpha_V=0.0005\ /^{\circ}\mathrm{C}$。求膨胀水箱的最小容积 $V_{\min}$。
\tuyi{习题 1.10 图：一个闭式水暖系统示意图，自下而上依次为\textbf{锅炉}$\to$\textbf{管路}$\to$\textbf{散热器}$\to$回到锅炉，管路最高处（水暖系统顶部）接一个\textbf{膨胀水箱}；系统内水的总体积为 $10\ \mathrm{m^3}$，膨胀水箱用来容纳水温升高后多出来的那一部分水，使管路不被胀裂。}
\xstarfull{2}\yiju{教材式 1.8（体积膨胀系数的定义，E6）；E2"第 1--3 章偏概念、掌握基本公式与解释即可"；E4 历年真题中未见独立成题，故压至两星}\nandu{易}

\item 相距 $d=2\ \mathrm{mm}$ 的两块平板插入水中，水的表面张力系数 $\sigma=0.0725\ \mathrm{N/m}$，接触角 $\theta=8^{\circ}$，求两板间的毛细水柱高 $h$。
\tuyi{习题 1.11 图：两块竖直平行平板插入水中，两板\textbf{内间距} $d=2\ \mathrm{mm}$；两板之间水面升高成凹形液面（水能润湿板面），液面与板壁的\textbf{接触角}为 $\theta$；要求板间水柱高出外部自由液面的高度 $h$。}
\xstarfull{1}\yiju{E5 题库与历年真题中未见毛细现象成题；E6 教材 1.3.4 节，教材原文明言"由于表面张力很小，一般对液体的宏观运动不起作用，可以忽略不计"，仅在流体计量、液滴与气泡形成等问题中才必须考虑}\nandu{中}

\end{ti}
